% !TEX TS-program = pdflatex
% !TEX encoding = UTF-8 Unicode

\begin{otherlanguage}{portuguese}

\chapter*{Agradecimentos}
\markboth{Agradecimentos}{}

A conclusão desta dissertação não teria sido possível sem o apoio e o incentivo de muitas pessoas, às quais manifesto o meu mais sincero agradecimento.

Em primeiro lugar, agradeço aos meus orientadores, Professora Maryam Abbasi e Professor Pedro Martins, sem cuja orientação este trabalho não teria chegado à sua forma atual. Esta dissertação nasceu de uma proposta apresentada pela Professora Maryam e foi sendo desenvolvida através do acompanhamento contínuo de ambos. Agradeço-lhes a disponibilidade, os comentários e sugestões construtivas que partilharam em todas as etapas do processo de investigação e a iniciativa de disponibilizarem materiais de orientação e recursos que ajudaram a definir tanto a metodologia como o trabalho escrito.

Agradeço de forma particular aos Professores Filipe Madeira e Filipe Cardoso, que, ao longo dos anos, partilharam generosamente a sua experiência profissional e me deram incentivos que levarei comigo para a vida profissional. Para além desse apoio contínuo, quero recordar um momento específico que alterou o rumo dos meus estudos. Não me inscrevi na primeira fase de candidaturas porque tinha dúvidas sobre continuar. Procuraram-me durante uma aula e, em pouco tempo, conseguiram motivar-me e esclarecer-me quanto ao que poderia perder se desistisse. Essa conversa foi o ponto de viragem: inscrevi-me na segunda fase nesse mesmo dia. Graças a essa intervenção, consegui ainda entrar no Mestrado e tornar esta dissertação possível; por isso, estou-lhes profundamente grato.

Deixo um agradecimento muito especial aos meus amigos e colegas Vasco Geada e Ivo Camacho, que me têm acompanhado em todas as etapas do meu percurso académico no Instituto Politécnico de Santarém desde 2021. Desde o Curso Técnico Superior Profissional em Tecnologias e Programação de Sistemas de Informação, passando pela Licenciatura em Informática, até à frequência simultânea da Licenciatura em Engenharia Informática e deste Mestrado em Informática Aplicada (MIA), a sua amizade, motivação e camaradagem foram uma constante. Percorrer este caminho, repartido por quatro formações, não foi fácil, e ter ao meu lado pessoas que partilham o mesmo compromisso fez toda a diferença. Posso afirmar, com plena convicção, que a sua presença moldou as minhas ambições académicas de formas que não antecipava quando iniciei este percurso.

\end{otherlanguage}
